\documentclass[UTF8,12pt,a4paper,fontset=fandol]{ctexart}

% 支持从项目根目录、深度学习目录和当前专题目录编译。
\makeatletter
\def\input@path{{../}{../../}}
\makeatother
\usepackage{researchnotes}

\title{大语言模型中的键值缓存}
\author{}
\date{}

\begin{document}
\maketitle
\tableofcontents

\section{研究对象、直观动机与阅读路线}
\label{sec:kv-introduction}

\keyterm{大语言模型（large language model, LLM）}生成一段文字时，通常不是一次给出整段答案，而是在已有文本之后依次选择下一个\keyterm{词元（token）}。词元是分词器使用的离散单位，可能对应一个汉字、一个英文单词的一部分、标点或特殊控制符，不能简单等同于字数。假设提示包含三个词元 $x_1,x_2,x_3$，模型先根据这三个词元预测 $x_4$，再根据前四个词元预测 $x_5$。最直接的实现每次都把不断增长的全文重新送入网络；然而，在因果注意力中，前三个位置不允许查看后来的第四个位置，因此它们先前算出的许多结果没有改变。把这些结果保留下来，就能把下一次计算集中到新增位置上。

\keyterm{键值缓存（key--value cache, KV-cache）}是在自回归推理中，按层保存历史位置的键向量和值向量，以便后续查询直接使用的运行时状态。键用于计算相关性，值用于构造注意力的加权输出。它保存的不是自然语言答案，也不是模型参数，更不是一个简单的“词元编号到词向量”的字典。高层的键和值已经包含该词元此前的上下文，同一个词元出现在不同前缀中，通常会产生不同的缓存。理解这一点，才能正确讨论跨请求复用、修改历史消息和显存估算。

一个具体的工作过程是：模型处理“今天”“天气”“很好”三个示意词元后，每个注意力层都留下三组键值，并从最后位置的输出预测下一词元“，”。下一次前向计算只输入这个新词元；该位置在第一层产生新的查询、键和值，查询读取该层原有三组键值和当前一组键值，得到新的隐藏表示；此后依次通过更高层，每层都重复相同的增量操作。处理结束后，每层的缓存增长到四个位置，模型又获得预测第五个词元的分布。这里的分词仅为说明过程，不主张实际分词器一定采用这样的切分。

缓存带来的主要收益是避免反复计算历史位置的投影、注意力输出、前馈网络和层间表示；付出的主要代价是让历史键值持续占用内存。历史越长，新查询仍须访问的键值越多。因此，缓存同时连接三个问题：在数学上，哪些中间量保持不变；在算法上，哪些计算能够省略；在系统上，有限显存能同时维持多少条正在生成的序列。只说“缓存可以加速”不足以判断实际瓶颈，也无法解释为什么一个模型的参数装得下，长对话却仍会显存不足。

本文以参数固定、关闭训练随机性、采用因果自注意力的仅解码器 Transformer 为基本对象。先建立注意力与缓存的数学记号，证明增量计算与完整重算的等价条件，再推导时间与存储复杂度，随后介绍位置编码、掩码、实现接口和服务系统。多头结构、分页、前缀共享、量化与淘汰分别改变不同部分，不能混为同一类优化。最后给出实验协议、排错方法和带解答的练习。基础结构参见 Transformer 原始论文\cite{transformer}，软件缓存接口的概念可参见官方说明\cite{hf-cache}；本文的复杂度与容量算例均在文中给定的假设下推导，不代表某一商用模型的实测结果。

\section{从自回归分解到注意力张量}
\label{sec:kv-attention}

\subsection{预测目标与统一记号}

设词表为有限集合 $\mathcal V$，随机词元序列为 $X_1,\ldots,X_T$，其一次实现写为 $x_1,\ldots,x_T$；$\theta$ 表示固定的模型参数。自回归语言模型使用条件分解
\begin{equation}
 p_\theta(x_1,\ldots,x_T)
 =\prod_{t=1}^{T}p_\theta(x_t\mid x_1,\ldots,x_{t-1}).
 \label{eq:kv-autoregressive}
\end{equation}
第一个条件分布可以由专门的序列起始词元定义。模型在位置 $t$ 读入 $x_t$ 后形成分数向量 $\mathbf z_t\in\mathbb R^{|\mathcal V|}$，经 softmax 得到对 $X_{t+1}$ 的分布。\keyterm{分数（logits）}是归一化之前的实数，不能直接当作概率。采样器根据分数和温度等配置选择新词元；缓存服务于模型前向计算，与采用贪心、随机采样还是束搜索是不同层面的事情。

全文用 $L$ 表示注意力层数，$d$ 表示模型隐藏维度，$H_q$ 表示查询头数，$H_{kv}$ 表示键值头数，$d_k,d_v$ 分别表示单头键和单头值的维度。常见结构令 $d_k=d_v=d_h$ 且 $d=H_qd_h$，但这些相等关系不是所有架构的必然性质。$B$ 表示当前批中序列数，$P$ 表示提示长度，$G$ 表示计划输出的词元数，$T$ 表示某次前向完成后已处理的词元数。$b$ 表示一个缓存元素的存储字节数；它与批量大小 $B$ 大小写不同。若序列长度不等，则第 $i$ 条序列的有效长度记为 $T_i$。

为与程序中的矩阵乘法一致，本文把单位置表示写为行向量。第 $\ell$ 层的输入表示为 $\mathbf h_t^{(\ell-1)}\in\mathbb R^{1\times d}$；其具体输入可能先经归一化。下文用 $\mathbf u_t^{(\ell)}\in\mathbb R^{1\times d}$ 统称送入该层注意力投影的表示，避免把前置归一化与后置归一化结构混同。单头情形的投影为
\begin{equation}
 \mathbf q_t=\mathbf u_tW_Q,\qquad
 \mathbf k_t=\mathbf u_tW_K,\qquad
 \mathbf v_t=\mathbf u_tW_V,
 \label{eq:kv-projection}
\end{equation}
其中 $W_Q,W_K\in\mathbb R^{d\times d_k}$、$W_V\in\mathbb R^{d\times d_v}$；暂时省略层与头索引，以及可能存在的偏置和位置变换，后文会说明位置变换发生在哪里。

\subsection{因果注意力与三个投影的分工}

把前 $T$ 个位置的键、值按行堆叠成 $K_{1:T}\in\mathbb R^{T\times d_k}$ 和 $V_{1:T}\in\mathbb R^{T\times d_v}$。位置 $t$ 的注意力输出是
\begin{equation}
 a_{tj}=\frac{\exp(s_{tj})}{\sum_{r=1}^{t}\exp(s_{tr})},\quad
 s_{tj}=\frac{\mathbf q_t\mathbf k_j^{\mathsf T}}{\sqrt{d_k}},\quad
 \mathbf o_t=\sum_{j=1}^{t}a_{tj}\mathbf v_j.
 \label{eq:kv-attention-row}
\end{equation}
查询 $\mathbf q_t$ 描述当前位置要寻找什么，键 $\mathbf k_j$ 提供位置 $j$ 如何被匹配的表示，值 $\mathbf v_j$ 提供匹配后传递的内容。三者都是训练得到的数值表示，这种语言只是解释其计算角色，不意味着每一维都有可命名的语义。分母必须在当前允许访问的所有键上共同归一化，因此增加一个可见键会改变新查询对应的整行权重。

一次处理整个序列时，写成矩阵形式为
\begin{equation}
 O=\operatorname{softmax}_{\mathrm{row}}
 \left(\frac{QK^{\mathsf T}}{\sqrt{d_k}}+M\right)V,
 \quad M_{ij}=\begin{cases}0,&j\leq i,\\-\infty,&j>i.\end{cases}
 \label{eq:kv-masked-attention}
\end{equation}
softmax 按行作用，$M$ 是\keyterm{因果掩码（causal mask）}。实际程序也可能用布尔掩码，或使用内核中的隐式因果规则；这些是表示形式的区别。禁止读取未来位置的条件才是本质。训练时所有输入位置已知，因此可以并行计算整个下三角区域；生成时下一个输入尚未选出，序列方向上存在无法由缓存消除的依赖。

多头注意力为不同头使用不同的投影，在每个头内计算式\eqref{eq:kv-attention-row}，再连接输出并经过输出投影。对于一个批次，常见缓存逻辑形状是 $[B,H_{kv},T,d_h]$，分别对应批、键值头、序列、头内维度。具体内核可能交换轴、按块存储或将维度分块，因而不能仅凭物理排列猜测语义。缓存形状中的头数是实际存储的 $H_{kv}$，而计算注意力输出所需的查询头数是 $H_q$；两者的区别将在第\ref{sec:kv-heads}节直接影响显存估算。

\subsection{一个可手算的增量注意力例子}

\begin{example}[相同历史键值面对新的查询]
考虑单层单头，令 $d_k=d_v=1$。历史缓存为 $K_{1:2}=(0,\log 2)^{\mathsf T}$、$V_{1:2}=(2,5)^{\mathsf T}$。新位置产生 $q_3=1$、$k_3=\log 3$、$v_3=8$，求位置三的注意力输出。
\end{example}
\begin{solve}
当前可见的分数为 $(0,\log 2,\log 3)$，指数为 $(1,2,3)$，归一化后权重为 $(1/6,2/6,3/6)$，故 $o_3=2/6+10/6+24/6=6$。旧的 $k_1,k_2,v_1,v_2$ 可直接读取；新的三个权重仍须由 $q_3$ 计算，不能沿用上一查询的权重。若另一个查询为 $q=0$，对同样三个键的权重变为均匀分布，输出为 $5$。缓存中相同的键值并不意味着每一步的注意力分布相同。
\end{solve}

这个例子同时说明为什么通常缓存键和值而不缓存所有历史查询。后续位置要访问的是“历史有哪些可匹配内容”，因此需要历史 $K,V$；历史查询仅用于计算历史位置自身的输出，该输出已经向更高层传递完毕。对于标准因果增量生成，后续位置不再调用历史查询。实现为了调试、可视化或特殊算法保存查询是另外的需求，不能据此把普通 KV-cache 的容量公式乘上三个投影。

\section{缓存复用为何正确：前缀不变性}
\label{sec:kv-correctness}

\subsection{需要明确的假设}

“旧词元不变，所以旧键值不变”只有补全模型条件才成为严格论证。旧词元的输入表示必须不因追加文本而改变；旧位置的每一层计算必须只依赖允许访问的历史；参数、归一化规则、位置编码规则和外部条件必须固定。常见的 LayerNorm、RMSNorm 在单个位置的特征维上归一化，不跨序列长度重新统计，因此符合这个要求；若某个自定义层用全序列均值归一化所有位置，就未必符合。推理时还应关闭 dropout 等训练随机性，或严格保持对应随机状态一致。

\begin{proposition}[前缀不变性]
\label{prop:kv-prefix-invariant}
设一个 $L$ 层网络的输入嵌入及位置表示对旧位置保持固定；每层在位置 $j$ 的输出仅依赖前一层位置 $1,\ldots,j$ 的表示和固定参数，各操作均为确定性函数。对任意长度 $T$ 的输入及其长度 $T+r$ 的扩展，旧位置 $j\leq T$ 在每一层的表示相同，因而该位置由这些表示产生的键和值也相同。
\end{proposition}
\begin{proof}
在第零层，两次输入的前 $T$ 个词元及其位置表示相同，所以相应嵌入相同。假设第 $\ell-1$ 层所有旧位置表示相同。对任一 $j\leq T$，第 $\ell$ 层在位置 $j$ 读取的仅是位置 $1,\ldots,j$ 的前层表示，它们由归纳假设完全相同；参数和计算规则也相同，因此确定性运算的输出相同。对层数归纳便得到全部旧位置的各层表示相同。键和值是这些表示在固定投影和固定位置规则下的函数，所以同样保持不变。
\end{proof}

缓存正确性还需要一个跨时间的不变式：处理完长度 $T$ 的前缀后，每层缓存应恰好对应完整前向计算得到的前 $T$ 个键值，且逻辑顺序、位置编号和有效掩码一致。若这一不变式成立，计算新位置时读取的历史键值与完整重算相同，当前键值和查询又由相同的新位置输入逐层产生，因此新位置每一层的输出相同。将新键值追加后，不变式扩展到 $T+1$。这给出了增量解码的归纳证明，而不是仅凭最终生成文本相似来判断实现正确。

\subsection{等价性的适用范围}

上述结论是在精确算术下的函数等价。浮点计算会受到求和次序、矩阵乘法内核、精度和批量形状影响，因此完整重算与缓存版本未必逐位一致。合理的验证方式是在固定输入轨迹上比较分数向量或中间表示，采用与精度相适应的绝对、相对容差。若两个候选词元的分数非常接近，小的数值差也可能改变贪心选择，随后两个生成轨迹就会分叉；不能把轨迹分叉一概视为缓存逻辑错误，也不能只看首个词元相同就宣布数值等价。

\begin{example}[双向注意力不能直接沿用同样的证明]
设某层没有因果掩码，两个位置的分数都为零，值分别为 $0$ 和 $2$，那么位置一的输出为 $1$。追加值为 $10$ 且分数仍为零的第三个位置后，位置一的输出变为 $4$。旧位置输入词元没有变化，但它可以读取新增位置，故旧输出改变；更高层由旧输出产生的键值也会改变。
\end{example}

相同问题会出现在修改前缀内部词元、改变模型权重、切换影响投影的适配器，或者改变旧位置所用位置变换时。一般情况下，修改第 $j$ 个词元后，位置 $j$ 及之后的缓存应视为失效；位置 $j$ 之前仍满足不变条件的部分可以复用。这是保守而可验证的规则。个别架构可能允许更细粒度复用，但应依据实际依赖图证明，不能仅比较被修改词元本身。

缓存中的每层状态也不能相互替代。同一个位置在第一层和第十层的表示经过了不同变换，键值承担不同的寻址空间和内容空间；即使张量形状相同，它们也不是同一个数学对象。将层索引错位常常不会触发形状报错，却会破坏模型输出。正确实现必须把层、头、位置、请求身份和数据类型共同作为状态契约的一部分。

\section{预填充、解码与生成循环的时间顺序}
\label{sec:kv-phases}

\subsection{预填充完成了什么}

\keyterm{预填充（prefill）}是对初始提示进行前向计算并建立缓存的阶段。长度为 $P$ 的提示在每层产生 $P$ 个位置的查询、键和值，因果注意力允许不同位置的行并行计算。完成所有层之后，每层保存提示对应的键值，最后一个有效位置的分数给出第一个输出词元的分布。提示中的早期位置虽然不直接用于最后一步采样，却通过各层键值影响后续位置，不能简单跳过。只计算最后位置而没有任何历史中间状态，无法完成一般多层 Transformer 的完整提示处理。

预填充并不要求把全部注意力矩阵长期保存。注意力矩阵是当前前向的中间结果，缓存则是跨生成步骤持续存在的键值。优化内核可以边计算边归约注意力，减少中间矩阵的存储，最后仍留下需要的 $K,V$。因此，“没有显式存储 $P\times P$ 矩阵”不等于“没有 KV-cache”，也不等于“取消了稠密注意力中的位置两两交互”。

\subsection{为什么第一个输出不需要额外的一次解码}

\keyterm{解码（decode）}在这里特指提示之后的增量前向阶段，而“解码算法”有时也用来泛指从分数选择词元的策略。为了避免混淆，下面把选择词元称为采样或选择。预填充最后位置的分数已经能够选择 $x_{P+1}$，此时缓存长度仍为 $P$，因为新词元刚被选出，还没有作为输入通过各层。只有下一次把 $x_{P+1}$ 送入模型后，缓存才变为 $P+1$，并产生预测 $x_{P+2}$ 的分数。

因此，若预填充产生第一个输出，计划生成 $G\geq1$ 个词元，通常需要一次预填充和 $G-1$ 次单词元增量前向。结束时，最后一个刚选出的词元可能尚未进入缓存。某些接口会额外处理它以便延续会话，或者采用不同的返回状态约定；判断缓存长度时必须看接口究竟承诺“已经输出”还是“已经计算”。这一时间差是重复输入末尾词元、漏掉末尾词元及位置偏移的常见来源。

例如 $P=3,G=3$，预填充读取 $x_1,x_2,x_3$，选择 $x_4$，缓存长度为三；第一次增量前向读取 $x_4$，选择 $x_5$，缓存长度为四；第二次读取 $x_5$，选择 $x_6$，缓存长度为五。用户已看到六个词元的整段文本，而缓存只处理到第五个位置。如果希望继续预测第七个词元，需要先把第六个词元送入模型。这种差异不是缓存损坏，而是生成循环的自然时序。

\subsection{分块输入与终止条件}

预填充可以拆为多个块，每个块读取已建立的历史缓存，并在块内使用因果掩码。若历史已有 $T_0$ 个位置，新块长度为 $U$，则该块第 $i$ 个查询的全局位置是 $T_0+i$，它可见全部历史和块内前 $i$ 个位置。在位置编号、掩码和数值规则一致的条件下，分块与一次性预填充计算相同的函数。分块改变调度与临时内存需求，不自动改变最终需要保留的全局注意力缓存总量。

生成循环通常在选出结束词元、满足外部停止规则或达到最大输出长度时终止。不同请求会在不同时间结束，批量系统必须为每条请求维护独立的已处理长度和终止标记。终止后释放哪些缓存还取决于是否需要继续对话、是否保留公共前缀以及共享引用计数，不能直接清空整个批的缓存。缓存释放是资源生命周期问题，停止采样则是序列生成问题，两者相关但不等价。

\section{计算复杂度：省下的部分与保留的部分}
\label{sec:kv-complexity}

\subsection{单次前向的成本模型}

先考虑标准多头稠密注意力，设所有层宽度相同，前馈中间维度为 $d_{\mathrm{ff}}$，且忽略词表输出投影、softmax 的低阶项和常数。令
\begin{equation}
 C_{\mathrm{tok}}=d^2+d\,d_{\mathrm{ff}}
 \label{eq:kv-token-cost}
\end{equation}
表示单位置投影与前馈运算的量级。长度为 $n$ 的完整前向总成本为 $O(L(nC_{\mathrm{tok}}+n^2d))$。其中第一项来自逐位置线性层和前馈层，第二项来自查询键乘积以及对值的加权。该估计假定头维度之和与 $d$ 同阶；采用特殊头结构时，应回到实际投影形状计数。

缓存版本处理一个新位置时，历史位置的逐位置计算已经完成。新位置只支付 $O(LC_{\mathrm{tok}})$ 的新投影和前馈成本，同时在每层读取并匹配约 $n$ 个历史位置，支付 $O(Lnd)$ 的稠密注意力成本。因此单步是 $O(L(C_{\mathrm{tok}}+nd))$，而不是关于上下文长度的常数。这里的 $n$ 可以理解为包含当前词元的可见长度，差一项不会影响阶数，但实现中的索引必须精确。

\subsection{整段生成的求和}

固定提示长度 $P$，生成 $G\geq1$ 个词元，并采用第\ref{sec:kv-phases}节的时序。无缓存且每次完整重算时，前向输入长度依次是 $P,P+1,\ldots,P+G-1$，所以
\begin{equation}
 C_{\mathrm{no\ cache}}=
 O\left(L\sum_{j=0}^{G-1}
 \bigl((P+j)C_{\mathrm{tok}}+(P+j)^2d\bigr)\right).
 \label{eq:kv-nocache-cost}
\end{equation}
缓存版本只做一次完整提示处理，随后每次处理一个新位置，故
\begin{equation}
 C_{\mathrm{cache}}=O\left(L\left[
 (P+G-1)C_{\mathrm{tok}}+
 d\left(P^2+\sum_{j=1}^{G-1}(P+j)\right)\right]\right).
 \label{eq:kv-cache-cost}
\end{equation}
式中的 $P^2$ 代表预填充的稠密注意力量级，不是在断言因果区域恰有 $P^2$ 个有效元素；精确的有效查询键配对数为 $P(P+1)/2$。后续求和为 $(G-1)P+G(G-1)/2$。这些表达式说明，缓存既减少历史投影与前馈重算，也减少重复计算的历史注意力行。

当忽略提示、固定模型宽度并令生成长度增长时，每次从头重算的注意力累计量级为 $O(G^3)$，增量缓存为 $O(G^2)$。这句常见结论针对的是特定朴素基线和稠密注意力项，并不是说单步复杂度从三次变为二次，也不是说一次训练前向需要三次复杂度。若只需要提示后的第一个词元，两者均需完成预填充，缓存几乎没有机会省下后续重算；若提示很长且输出也很长，复用的收益则更加突出。

\begin{example}[统计实际处理了多少个位置]
提示长度 $P=1024$，生成 $G=128$ 个词元。朴素实现用于投影和前馈的位置总数为 $128\times1024+128\times127/2=139200$；缓存实现为 $1024+127=1151$。两者位置数比约为 $120.94$，但这个比值仅反映该部分算术工作，不能当作端到端速度提升。缓存版本仍需计算新查询的历史注意力，并支付内存读取、内核启动与采样成本。
\end{example}

\subsection{算术量不直接等于延迟}

图形处理器（graphics processing unit, GPU）的计算单元适合成批矩阵运算；一次只处理一个位置时，权重读取与启动开销可能相对突出。预填充通常具有更大的矩阵维度和较高的数据复用机会，单请求解码则更容易受到内存带宽限制，但这只是需要验证的硬件工作区间，而非所有设备、模型和批量下的普遍定理。增大批量有时能让多个请求共享一次权重读取，却同时增加各请求独立的缓存访问。

把每步的浮点运算量记为 $F$，从主要存储层搬运的字节数记为 $D$，设备对应的可用运算吞吐和带宽分别记为 $\Pi$ 与 $\beta$。忽略重叠细节时，一个有用的理想下界是 $\max(F/\Pi,D/\beta)$。实际耗时还含调度、通信和同步，因此不应由峰值指标直接预测性能。缓存显著降低 $F$，却仍需读取历史 $K,V$；当 $D/\beta$ 成为主导，进一步减少乘法可能收效有限，缩小缓存或改善访存才更有针对性。

\section{显存占用与容量规划}
\label{sec:kv-memory}

\subsection{从元素个数推导字节数}

对每条序列、每个位置、每层和每个键值头，需保存一个 $d_k$ 维键和一个 $d_v$ 维值。若所有层结构相同、键值精度相同且批内长度都为 $T$，纯缓存载荷为
\begin{equation}
 M_{KV}=LBT H_{kv}(d_k+d_v)b.
 \label{eq:kv-memory-general}
\end{equation}
常见的等维情形变为
\begin{equation}
 M_{KV}=2LBT H_{kv}d_hb,
 \qquad m_{\mathrm{token}}=2LH_{kv}d_hb,
 \label{eq:kv-memory-equal}
\end{equation}
其中 $m_{\mathrm{token}}$ 是一条序列增加一个已处理词元时，跨全部层增加的缓存字节数。前面的因子二来自键和值各一份，不来自训练梯度或两个批次。如果键和值的精度不同，则把 $d_kb_k+d_vb_v$ 代入式\eqref{eq:kv-memory-general}中的相应项。

对于无填充浪费、没有共享的变长批次，应把 $BT$ 换成 $\sum_{i=1}^{B}T_i$。如果各层的键值头数、维度或保留窗口不同，更一般的表达为
\begin{equation}
 M_{KV}=\sum_{i=1}^{B}\sum_{\ell=1}^{L}
 T_{i,\ell}H_{kv,\ell}
 (d_{k,\ell}b_{k,\ell}+d_{v,\ell}b_{v,\ell}).
 \label{eq:kv-memory-heterogeneous}
\end{equation}
这里 $T_{i,\ell}$ 是第 $i$ 条序列在第 $\ell$ 层实际保留的有效位置数。此式仍只计算有效载荷；块对齐、预留空间、索引、量化尺度、共享计数和分配器碎片需要另外计入。先算数学上必需的元素数，再解释实现多分配了什么，比直接观察一个显存读数更容易定位问题。

\subsection{完整的数值算例}

\begin{example}[同一宽度下的缓存规模]
设 $L=32,H_q=32,d_h=128$，所有缓存使用每元素两字节的数据类型。分别考虑 $H_{kv}=32$ 的多头结构和 $H_{kv}=8$ 的分组结构，求单序列在 $T=8192$ 时的缓存载荷。
\end{example}
\begin{solve}
前者每词元占 $2\times32\times32\times128\times2=524288$ 字节，即 $512\,\mathrm{KiB}$；总量为 $4294967296$ 字节，即 $4\,\mathrm{GiB}$。后者每词元占 $131072$ 字节，即 $128\,\mathrm{KiB}$，总量为 $1\,\mathrm{GiB}$。若同时有八条等长序列且不共享缓存，总量分别为 $32\,\mathrm{GiB}$ 与 $8\,\mathrm{GiB}$。这里 $1\,\mathrm{KiB}=2^{10}$ 字节、$1\,\mathrm{GiB}=2^{30}$ 字节；厂商或监控中使用十进制 GB 时，数值需要转换。
\end{solve}

这个算例没有绑定某个现成模型，避免把不同架构配置误套在一起。两字节可以对应不同浮点格式，容量相同不意味着数值性质相同。类似地，模型权重采用四比特量化，不能推出缓存也是四比特：权重是常驻参数，缓存是输入驱动的中间状态，二者的精度通常由不同机制控制。估算时必须从实际模型配置读取键值头数，并确认缓存的数据类型与保留策略。

假设设备总显存为 $24\,\mathrm{GiB}$，权重、工作区和保留余量共占 $16\,\mathrm{GiB}$，可供缓存使用的预算为 $8\,\mathrm{GiB}$。沿用每词元 $128\,\mathrm{KiB}$ 的结构，理论上能容纳 $65536$ 个互不共享的有效词元。若每请求当前缓存为 $8192$ 个词元，载荷上限对应八条请求；但若还要生成更长输出，或者工作区随批量增长，这个上限就不能直接用作安全并发数。容量规划要说明预算是当前时刻、请求峰值还是系统允许的长期上限。

\subsection{为什么参数装得下仍会内存不足}

推理的总显存可概念性拆成权重、KV-cache、当前激活与临时工作区、运行时及通信缓冲四部分。权重在固定模型下相对稳定，缓存随活跃序列长度与并发增长，临时工作区则可能在长提示预填充或大批量时突然升高。内存分配器还可能保留已经申请但暂时未使用的块，所以“已分配”“已保留”和设备工具报告的占用不必相同。测量应说明指标口径，不能把它们相减后直接认定为缓存泄漏。

请求长度分布也比平均长度更重要。两个系统可能平均各有四条请求，但一个系统总是处理短输入，另一个系统会偶尔同时接收多个长文档；后者的缓存峰值和排队风险完全不同。服务端常设置总缓存词元预算、单请求最大长度和预填充预算，并在接纳请求时判断剩余空间。预留全部最大输出较保守，边生成边申请较灵活，但需要定义容量不足时暂停、换出或重算的策略；不存在完全不考虑未来增长又保证永不超额的免费方案。

\section{多头、多查询、分组查询与潜在表示}
\label{sec:kv-heads}

\subsection{查询头数与键值头数为何不同}

\keyterm{多头注意力（multi-head attention, MHA）}通常为每个查询头配备独立的键和值，因此 $H_{kv}=H_q$。\keyterm{多查询注意力（multi-query attention, MQA）}让所有查询头共享一组键和值，即 $H_{kv}=1$，而查询头仍可以有多个\cite{mqa}。\keyterm{分组查询注意力（grouped-query attention, GQA）}位于两者之间，让一组查询头共享一个键值头\cite{gqa}。这些是模型的投影组织方式，必须与训练所得参数一致，不能在普通多头模型上随意删除键值头就声称是等价推理。

设 $H_q$ 可被 $H_{kv}$ 整除，每组有 $g=H_q/H_{kv}$ 个查询头。令 $\gamma(h)=1+\lfloor(h-1)/g\rfloor$ 表示查询头 $h\in\{1,\ldots,H_q\}$ 所属的键值组，则该头输出为
\begin{equation}
 O_h=\operatorname{softmax}_{\mathrm{row}}
 \left(\frac{Q_hK_{\gamma(h)}^{\mathsf T}}{\sqrt{d_h}}+M\right)
 V_{\gamma(h)}.
 \label{eq:kv-gqa}
\end{equation}
虽然同一组共享 $K,V$，不同查询头的 $Q_h$ 不同，注意力权重仍可以不同。因此共享键值不等于共享所有注意力输出，也不等于把查询头数减少为一。

在层数、长度、头维度和精度相同的条件下，GQA 相对 MHA 的缓存载荷比例为 $H_{kv}/H_q$；MQA 则为 $1/H_q$。这是元素计数的直接结果。注意力打分与值聚合仍要为各查询头进行，所以算术量并不必然按相同比例下降；读取共享键值时是否真正复用数据也取决于内核。若实现先将共享的键值物理复制成 $H_q$ 份再长期缓存，就会抵消原本的容量优势。逻辑广播、短暂展开和永久重复存储应在代码审查中分开判断。

减少键值头会约束模型可表达的投影结构，质量与效率之间需要训练和评估。GQA 论文研究了把已有多头检查点转换并继续训练的方式，而不是证明所有模型都能零训练、零损失地转换\cite{gqa}。阅读性能结果时，应区分架构本身的容量比例、特定内核的速度和特定训练配置的任务质量；这三个结论不能互相替代。

\subsection{潜在表示压缩的代数直觉}

另一种方向不是减少键值头，而是利用较低维的共享潜在表示生成键和值。为说明代数结构，设行向量 $\mathbf c_j\in\mathbb R^{1\times r}$，并令 $\mathbf k_j=\mathbf c_jU_K$、$\mathbf v_j=\mathbf c_jU_V$，其中 $U_K\in\mathbb R^{r\times d_k}$、$U_V\in\mathbb R^{r\times d_v}$。则
\begin{equation}
 \mathbf q_t\mathbf k_j^{\mathsf T}
 =(\mathbf q_tU_K^{\mathsf T})\mathbf c_j^{\mathsf T},
 \qquad
 \sum_j a_{tj}\mathbf v_j
 =\left(\sum_j a_{tj}\mathbf c_j\right)U_V.
 \label{eq:kv-latent-identity}
\end{equation}
若模型确实具有这样的参数化，计算可以围绕 $\mathbf c_j$ 进行，而不必长期保存展开后的全部键值。这个等式是矩阵乘法结合律的应用，不是对任意现成键值做低秩近似后仍然精确的保证。若低秩关系仅近似成立，就另外引入压缩误差。

DeepSeek-V2 中的\keyterm{多头潜在注意力（Multi-head Latent Attention, MLA）}采用键值联合压缩，并为位置相关部分设计解耦处理\cite{mla}。这个例子提示：真实缓存可能包含潜在向量以及额外的位置相关键，不能一律套用 $2LH_{kv}d_hb$。式\eqref{eq:kv-latent-identity}只是理解思路的简化模型，未复述其完整架构；在真实实现中，哪些矩阵可以吸收、位置变换如何组合以及临时展开占多少空间，仍应按模型定义和内核核验。

从优化层面看，GQA 改变键值头轴，潜在表示改变保存对象的表示维度，量化改变每元素的字节数，淘汰改变保留的位置数，分页则改变分配方式。它们对应容量公式的不同因素，因此可以有组合空间，也会产生接口约束。先指出正在缩小哪个因素，比把所有方法都称为“压缩 KV”更有助于分析质量代价和工程收益。

\section{位置编码、掩码与变长批次}
\label{sec:kv-position-mask}

\subsection{缓存位置不是张量下标的同义词}

位置编码使模型区分相同内容在不同顺序中的作用。对于将绝对位置向量加到嵌入上的结构，缓存中已经间接包含相应位置的信息；继续生成必须使用新词元的正确位置。对于\keyterm{旋转位置编码（rotary position embedding, RoPE）}，位置变换通常施加在查询和键上，而不是普通值向量上，其基本构造见文献\cite{rope}。工程上常直接缓存变换后的键，这样旧键无需每步重新旋转；也可以保存其他形式，但必须在读取路径上保持等价的变换约定。

为给出一个明确的约定，设 $R_p$ 是位置 $p$ 的正交旋转矩阵，未旋转的行向量为 $\bar{\mathbf q}_p,\bar{\mathbf k}_p$，定义 $\mathbf q_p=\bar{\mathbf q}_pR_p$、$\mathbf k_p=\bar{\mathbf k}_pR_p$。于是 $\mathbf q_t\mathbf k_j^{\mathsf T}=\bar{\mathbf q}_tR_tR_j^{\mathsf T}\bar{\mathbf k}_j^{\mathsf T}$。在固定频率的旋转组构造中，旋转组合体现相对位置关系；但这不意味着程序可以任意更改已缓存键的位置而不补偿。若新查询错误地使用位置零，旧键仍按原位置旋转，两者的相对关系就被破坏。

尤其需要避免两种操作：把已经旋转过的旧键再次旋转，以及只移动缓存张量却把物理下标当作新的逻辑位置。滑动窗口和环形缓冲区可能反复使用同一物理槽位，槽位编号与原始序列位置不再相同。实现必须维护从逻辑位置到存储槽的映射，并按模型要求决定是否保留绝对位置、采用相对偏置或进行经过证明的位置重映射。仅仅将张量裁短，不能证明位置语义仍然正确。

一些位置扩展方案可能根据当前总长度调整频率或缩放规则。如果追加序列会改变旧位置本应采用的变换，命题\ref{prop:kv-prefix-invariant}的假设就需要重新检查。固定同一套配置通常是缓存兼容的基础；动态规则则可能要求预先固定目标长度、重建相关状态或采用模型专门定义的修正，不能由“使用 RoPE”四个字自动推出可直接复用。

\subsection{矩形因果掩码的正确对齐}

已有历史长度 $T_0$，一次新输入 $U$ 个位置时，查询轴长度为 $U$，键轴长度为 $T_0+U$，注意力分数矩阵是矩形。按从一开始的编号，新块第 $i$ 行允许读取的键满足
\begin{equation}
 j\leq T_0+i,\quad 1\leq i\leq U,\quad 1\leq j\leq T_0+U.
 \label{eq:kv-rectangular-mask}
\end{equation}
若直接对这个矩形矩阵使用 $j\leq i$ 的左上角三角规则，就会把大量历史位置错误屏蔽。例如 $T_0=3,U=2$，第一行应可见位置一至四，第二行应可见一至五；错误规则却只允许第一行看第一个键。张量尺寸完全合法，输出却已不是目标模型。

当 $U=1$ 且键轴只包含有效历史与当前词元时，当前查询可以读取键轴全部位置，不再需要额外屏蔽“未来”。但是，如果静态缓存预分配了更长的容量，未写入的槽位仍须屏蔽；如果批中有填充位置，也要屏蔽填充。因而“单词元解码无需因果三角形”不能简化成“单词元解码不需要任何掩码”。对框架提供的因果标志，应核对其在非方形矩阵下的具体对齐语义；本文伪代码将显式构造允许关系，避免依赖版本相关的隐式约定。

\subsection{填充、位置编号与请求重排}

批处理中，短序列常用\keyterm{填充（padding）}补到共同长度。填充位置不能作为真实内容参加注意力归一化；左填充还会使张量列号偏离语义位置。常见做法是由有效词元掩码的累计计数生成位置编号，但这不是所有模型接口的统一约定。需要分别确认位置编号、缓存写入位置和有效长度的含义，而不是把它们都设成张量最后一维的长度。

假设一个批包含有效长度三和五的两条序列，左填充后张量宽度为五。两条序列下一词元的语义位置在从零编号时分别为三和五，而共同物理追加槽可能都是第六列。如果系统使用紧凑变长布局，两个请求又可能写入不同缓存块。此时“一个统一的标量缓存长度”不足以描述整个批，需要逐请求长度和存储映射。若使用右填充，直接取分数张量的最后一列还可能取到填充位置，应按最后有效位置收集分数，或遵循框架明确支持的输入约定。

\keyterm{连续批处理（continuous batching）}允许请求在不同生成步进入或退出，因此同一个批行号不保证始终对应同一个用户请求。缓存管理器必须随请求重排、压缩和调度更新映射；不能只重排输入词元而保留旧行的键值。最危险的错误不是越界，而是将合法形状的另一请求缓存交给当前查询。验证时应让不同请求使用容易区分的固定词元序列，并与逐请求独立运行的分数对照。

\section{从数学定义到最小正确实现}
\label{sec:kv-implementation}

\subsection{缓存对象的接口契约}

一个教学实现至少需要维护每层的键张量、值张量和已写入长度；若有批内变长，还需逐请求有效长度、位置与掩码；若有分页或共享，则进一步增加块表和引用计数。缓存对象的更新应明确是原地修改还是返回新对象。原地修改可以减少复制，却容易在分支生成中意外污染父前缀；返回新对象的语义更简单，但若内部深复制历史，则可能产生很高的带宽和容量成本。接口设计应把逻辑独立性与物理共享分开表达。

下面的伪代码使用单请求、固定模型、前置归一化结构和从零开始的位置编号。\texttt{project}、\texttt{attend} 等名称是抽象操作，不是对某个软件版本的调用承诺；\texttt{attend} 接受显式掩码并按式\eqref{eq:kv-gqa}处理头共享。缓存只保存键和值，当前块的隐藏表示在层间传递。每次调用结束后，各层长度应一致，否则说明某层漏写或重复写入。

\begin{verbatim}
forward_chunk(ids, cache):
    old = cache.length
    U = len(ids)
    pos = arange(old, old + U)
    h = embed(ids, pos)
    for ell in range(L):
        u = norm_attention[ell](h)
        q, k, v = project[ell](u)
        q, k = apply_position_if_needed(q, k, pos)
        K = logical_append(cache.K[ell], k)
        V = logical_append(cache.V[ell], v)
        allowed = key_positions(K) <= pos[:, None]
        o = attend(q, K, V, allowed)
        h = h + output_projection[ell](o)
        h = h + ffn[ell](norm_ffn[ell](h))
        cache.K[ell], cache.V[ell] = K, V
    cache.length = old + U
    logits = lm_head(final_norm(h))
    return logits, cache
\end{verbatim}

其中 \texttt{embed} 和 \texttt{apply\_position\_if\_needed} 必须按实际模型分工：使用绝对嵌入时在前者加入位置，使用 RoPE 时通常由后者变换查询和键，不能因为伪代码同时列出两个接口就重复加入同一种位置编码。为简化叙述，代码没有写填充掩码、局部窗口和模型特有层；这些扩展都应通过显式接口加入，而不能在不变式之外隐藏修改。

\subsection{生成循环与不能重复输入的前缀}

初始化时缓存为空，对完整提示调用一次 \texttt{forward\_chunk}；以后只输入尚未写入缓存的新词元。若缓存已经包含提示，而调用者又把提示连同新词元全部传回一个“只接收未处理输入”的接口，历史就会被重复追加。相反，有些高级生成接口接受完整输入后自行截取未处理部分，调用者也不能再擅自截取两次。正确做法是阅读接口契约，并在低层实现中断言追加长度与已处理长度一致；官方缓存文档可用于理解这些概念\cite{hf-cache}。

\begin{verbatim}
cache = empty_cache()
logits, cache = forward_chunk(prompt_ids, cache)
output_ids = []
for step in range(max_new_tokens):
    token = choose(logits[-1])
    output_ids.append(token)
    if is_stop(token) or step + 1 == max_new_tokens:
        break
    logits, cache = forward_chunk([token], cache)
\end{verbatim}

这里假定 \texttt{max\_new\_tokens} 为正整数；若允许零，应该直接返回空输出，避免无意义的提示计算。返回结果同时包含输出与缓存时，还应注明末尾词元是否已处理。若调用方希望缓存和全部输出严格对齐，可以在明确需要继续会话时处理最后词元，或额外返回“待处理词元”；选择哪一种约定都可以，但不能含糊。

\subsection{追加的语义与复制成本}

\texttt{logical\_append} 表示逻辑序列增长，不要求每步重新分配并复制整个张量。若每一步都通过拼接创建新连续缓冲，单层历史长度从一增长到 $T$ 时，累计复制量与 $1+2+\cdots+T$ 成正比，可能重新引入二次量级的内存搬运。这个开销并不改变数学上的缓存算法，却会显著影响实际速度。静态预分配用写指针原地填入新槽位，分页分配用新的块扩展逻辑序列，都是避免频繁复制的方式。

静态预分配把最大容量与有效长度分开，有利于稳定张量形状，但必须屏蔽未写槽位，并为最大容量承担空间成本。动态分配节省短序列的预留空间，但可能增加扩容、碎片和图编译管理成本。官方缓存策略文档介绍了动态、静态、换出与量化等实现类别\cite{hf-strategies}；选择时应依据已安装版本和模型支持情况，而不把类名当作跨版本不变的运行保证。本文的实现目标是展示状态流与正确性条件，不把这些伪代码称作已经测得高吞吐的推理引擎。

\section{数值稳定性与高效注意力内核}
\label{sec:kv-numerics}

\subsection{稳定的归一化与掩码边界}

softmax 的指数可能溢出。对一行有效分数 $s_1,\ldots,s_n$，先令 $m=\max_j s_j$，再计算 $\exp(s_j-m)/\sum_r\exp(s_r-m)$，其精确数学结果不变，而指数的最大值为一。掩码应使禁止位置的权重为零；若一行全被屏蔽，则没有可归一化的有效分布，直接进行减最大值可能产生未定义数值。正确模型通常保证每个有效查询至少可以访问自身或一个合法历史位置，填充查询则应按明确约定跳过或处理，不能寄希望于异常数值自动消失。

缓存存储精度、乘法精度和归约累加精度也可以不同。较低精度保存键值能够减小容量，但较高精度累加只能减轻后续数值误差，无法恢复写入时已丢失的信息。测试时需要分别记录这些设置。一个模型在较短序列上没有明显差异，不代表长序列上的大量点积、归一化与逐步反馈也具有同样误差；误差界、单步分数距离与最终任务质量是三个不同层次的观察。

\subsection{为什么可以不保存完整注意力矩阵}

对一个查询，把有效键分成若干块。已处理部分维护最大分数 $m$、指数和 $z$、未归一化的加权向量 $\mathbf u$；新块计算对应的 $m_b,z_b,\mathbf u_b$，其中各块内部均减去自己的最大分数。合并公式是
\begin{equation}
 \begin{aligned}
 m'&=\max(m,m_b),\\
 z'&=e^{m-m'}z+e^{m_b-m'}z_b,\\
 \mathbf u'&=e^{m-m'}\mathbf u+e^{m_b-m'}\mathbf u_b.
 \end{aligned}
 \label{eq:kv-online-softmax}
\end{equation}
处理全部块后输出为 $\mathbf u'/z'$。证明只需把两部分指数都换算到相同基准 $m'$，再将求和合并。空部分应单独初始化，避免直接运算两个负无穷之差。这个恒等式说明，精确注意力的归一化可以由块级统计量组合，不必先将全部分数写回大容量内存。

FlashAttention 的核心贡献之一是围绕存储层次组织分块计算，降低注意力中间结果的读写\cite{flash}。它与 KV-cache 解决的重复不同：前者减少一次注意力运算内部的存储流量，后者减少连续生成步骤之间的历史重算。二者可以同时存在。对于稠密全局注意力，采用这类内核并不自动把所需查询键配对的算术量降为线性，也不使跨步历史键值消失。预填充内核与单词元解码内核还可能因矩阵形状不同而采用不同并行策略。

当 GQA 的多个查询头共享键值时，理想内核会利用这种共享减少重复加载；当上下文很长而批量很小时，也可能把同一请求的历史分成多个段并行处理，再按式\eqref{eq:kv-online-softmax}合并统计量。分段越多，额外归约和临时存储也越多，因此并行粒度需要测量。数学上的可分解性提供实现空间，但并不直接保证某个分块尺寸最快。

\section{分页缓存、共享块与服务调度}
\label{sec:kv-paging}

\subsection{逻辑连续与物理连续}

一个请求的词元在逻辑上依次编号，不代表对应缓存必须放在一整块连续物理地址中。\keyterm{分页注意力（PagedAttention）}及其服务系统设计使用块表管理离散的缓存块，让注意力内核按逻辑顺序读取它们\cite{paged}。分页解决的是分配、增长、碎片和共享问题；只要读取的键值、位置和掩码不变，注意力仍计算原来的函数。它不要求删除历史，也不天然降低每个有效词元的数值精度。

设每块容纳 $S$ 个词元，长度为 $T$ 的请求需要 $\lceil T/S\rceil$ 块，预留槽位为 $S\lceil T/S\rceil$。不考虑特殊布局时，单请求末块未使用的槽位数在零与 $S-1$ 之间。这种有界的块内浪费与给每条请求预留最大长度不同，但块表、分配器与内核仍有开销。小块可减少尾部浪费并提高共享粒度，大块则可能简化管理并改善连续读取；不能只依据浪费槽位最少就选定块大小。

\begin{example}[块大小与尾部浪费]
取 $S=16$，三条互不共享的序列长度为 $17,31,48$，分别需要两块、两块、三块，总容量为 $112$ 个槽位，有效词元数为 $96$，尾部浪费为 $16$。若每条请求都预留 $64$ 个槽位，则总容量为 $192$，未用槽位为 $96$。这只是分配方式的容量比较，不包含块表开销，也不能推出访问速度之比。
\end{example}

\subsection{共享、引用计数与写时复制}

如果多条请求拥有相同且兼容的前缀，它们可以共同指向只读历史块，随后各自追加不同后缀。共享块需要\keyterm{引用计数（reference counting）}或等价的生命周期管理：只有没有活跃请求继续引用时，才允许回收或放入可淘汰池。释放一条请求的状态不等于立即释放它曾访问的全部物理块，否则另一个共享请求会读到失效内存。

部分写满的尾块尤其需要谨慎。假设两个分支共享同一块的前十个位置，随后都要写第十一个位置。若直接原地写入，后写分支会覆盖先写分支，因此应分配独立后继空间，或采用\keyterm{写时复制（copy-on-write）}复制必要内容后再分叉。完整的只读前缀可持续共享，发生修改的部分独立保存。这里的“复制”与每步复制整个历史不同；它只在共享状态需要分化时发生。

物理地址也不应成为缓存内容的身份。请求结束后，某个地址可能分配给另一段前缀；若旧块表或哈希索引没有同步失效，即使地址仍合法，也会产生错误复用。可靠的管理器需要将内容身份、所有权、有效长度与物理槽分开，并让分配、释放和查询操作遵循一致的并发协议。这些错误往往在高并发或取消请求时才出现，单请求顺序测试不足以覆盖。

\subsection{连续批处理与分块预填充}

连续批处理在每轮调度中重新选择活跃请求，允许完成的请求退出、新请求加入。它能减少固定批次等待最慢请求的浪费，但缓存总量成为接纳和调度的重要约束。调度器不仅要数“有多少请求”，还要数它们当前占多少块、下一步会新增多少块，以及待处理提示需要多大工作区。两个请求数相同的批次，其缓存压力可能相差一个数量级。

长提示预填充和短请求逐词元解码争用设备时，整段预填充可能占用较长的连续执行时间。将预填充分成较小块，并与解码工作交错，可以改变等待和响应的分布；块过小又会增加启动和调度次数，甚至降低矩阵运算效率。应以首词元延迟、词元间延迟、吞吐与内存峰值共同评价，而不是只看设备利用率。分块预填充保持的是同一请求内的计算依赖，不允许任意把后面的块先算完。

缓存不足时，服务系统可以停止接纳新请求、暂停部分请求、把状态换出到其他存储，或丢弃可重算的状态后再恢复。暂停保留状态仍占显存，换出消耗传输带宽，重算消耗设备时间；它们都把资源压力转移到其他成本。高负载时的尾部延迟往往由这些策略决定，不能仅凭一次无竞争的单请求测速评估整个服务。

\section{前缀缓存与跨请求复用的边界}
\label{sec:kv-prefix-sharing}

\subsection{公共前缀必须由相同计算产生}

\keyterm{前缀缓存（prefix caching）}把一次请求已经完成的兼容前缀状态保留下来，供另一请求跳过相同部分的预填充。普通 KV-cache 是同一生成过程内部的跨步复用，前缀缓存把复用扩展到多个请求或会话轮次。若公共前缀长度为 $P_s$，后缀长度分别为 $R_1,\ldots,R_B$，理想共享载荷为
\begin{equation}
 M_{\mathrm{shared}}=m_{\mathrm{token}}
 \left(P_s+\sum_{i=1}^{B}R_i\right),
 \label{eq:kv-prefix-memory}
\end{equation}
相比独立保存节约 $(B-1)P_sm_{\mathrm{token}}$。实际共享通常受块边界、引用管理和缓存淘汰影响。

复用条件首先是模型实际接收的词元前缀相同，而不只是显示出来的字符串看起来相似。聊天模板、角色标记、空格、特殊词元和分词器变化都可能改变输入编号；多模态模型中，相同占位符也可能对应不同图像或音频表示。模型权重及影响计算的适配器、位置规则、掩码和外部条件还必须兼容。一个合理的缓存身份可以概念性写为“前缀内容与其计算配置”的指纹；仅用文本字符串作为键是不充分的。

官方 vLLM 前缀缓存设计将父块身份、块内词元以及适配器或多模态等额外信息纳入块身份构造\cite{vllm-prefix}。其意义在于体现依赖关系，而不是认为某种哈希算法可以替代语义校验。相同文本片段如果出现在不同历史之后，高层键值通常不同，因此一般不能把任意文档段落的缓存像积木一样拼接。可直接复用的对象通常是相同完整前缀中的对应状态；支持非前缀复用需要额外模型假设或重计算方案。

\subsection{命中后仍然需要做什么}

前缀命中省去的是相同部分的前向重算。新的后缀仍须访问共享前缀键值，后续输出仍要逐步计算其条件分布，所以首词元延迟可能下降，而每步对历史的读取不一定减少。若缓存驻留在另一设备或主机内存，还需比较加载成本与重新预填充成本。只报告命中率不足以说明性能，应同时记录命中的词元数、实际避免的计算和额外传输。

还存在一个常被忽略的边界：普通 KV-cache 只保存各层键值，未必保存前缀最后位置的最终隐藏表示或词表分数。如果整条提示恰好全部命中，而系统需要立即给出下一词元分布，仅有这些键值未必足以直接得到分数。实现可以额外缓存边界分数或隐藏状态，也可以保留最后一小段供重算。重算时必须裁到一致的前缀边界，不能在已有完整前缀后又重复追加最后词元。所谓“全命中零计算”要以真实缓存内容和接口协议为前提。

对多轮对话，下一轮序列通常在上一轮生成结果之后追加用户消息，这有利于复用未变的历史；但聊天模板可能插入或改写结束标记，应用也可能截断旧消息或更新系统指令。系统应比较真正送入模型的最长兼容前缀，而不是只根据会话编号认为全部历史可用。若用户修改了早期消息，应从失效位置重建后续状态，不能仅替换显示文本而继续使用旧缓存。

\subsection{隔离与资源价值}

跨用户共享必须服从应用的数据隔离规则。缓存是输入驱动的中间表示，不能因为不是可读文本就当作没有数据含义；同时，命中与未命中的时延差异也可能暴露某段前缀是否曾被使用。工程上可把租户或安全域纳入复用命名空间，并明确过期、删除和监控策略。这里的目标不是禁止共享，而是让共享范围与应用实际允许的信息边界一致。

保留公共前缀也有机会成本：它占用的块本可服务其他活跃请求。长前缀但极少再次访问，未必比短而频繁访问的前缀更值得驻留。合理策略需要考虑未来复用概率、重算成本、加载成本、块大小和当前内存压力。最近最少使用等策略是可执行的启发式，并非所有工作负载的最优解。评估应给出冷热启动和复用间隔，防止把专门预热后的表现当作一般请求的默认性能。

\section{缓存量化及其误差来源}
\label{sec:kv-quantization}

\subsection{容量收益与实际开销}

\keyterm{量化（quantization）}用较少位宽表示缓存数值，通常同时保存尺度与零点。以一组数的仿射量化为例，取正尺度 $s$、整数零点 $z$ 和 $q$ 位整数范围 $[0,2^q-1]$，编码与重建为
\begin{equation}
 u=\operatorname{clip}\bigl(\operatorname{round}(x/s)+z,0,2^q-1\bigr),
 \qquad \widehat x=s(u-z).
 \label{eq:kv-quantization}
\end{equation}
没有截断且采用就近舍入时，单元素绝对误差不超过 $s/2$；存在截断时则需要另外考虑溢出范围的误差。较小尺度提高局部分辨率，却减少可覆盖范围，异常值与大动态范围因此影响量化设计。

若一共有 $N$ 个缓存标量，理想整数载荷是 $Nq/8$ 字节，但实际还要加尺度、零点、分组对齐和可能保留的高精度尾部。假设每 $g_q$ 个元素共享 $a$ 字节元数据，则仅元数据平均就增加 $a/g_q$ 字节每元素。由十六比特降为四比特，原始数值载荷缩为四分之一，并不保证总缓存、峰值显存或延迟也恰好缩为四分之一；临时反量化和模型权重都可能占据不可忽略的部分。

\subsection{键误差和值误差的不同传播路径}

设键和值重建误差分别为 $E_K,E_V$，原始分数为 $\mathbf s=\mathbf qK^{\mathsf T}/\sqrt{d_k}$。键误差引起的分数变化是 $\Delta\mathbf s=\mathbf qE_K^{\mathsf T}/\sqrt{d_k}$；对单个位置，由柯西不等式有 $|\Delta s_j|\leq\|\mathbf q\|_2\|\mathbf e_{K,j}\|_2/\sqrt{d_k}$。令原权重为 $\mathbf a$、变化后的权重为 $\widehat{\mathbf a}$，则输出差满足精确分解
\begin{equation}
 \widehat{\mathbf o}-\mathbf o
 =(\widehat{\mathbf a}-\mathbf a)V
 +\widehat{\mathbf a}E_V.
 \label{eq:kv-quant-error}
\end{equation}
键误差先改变寻址权重，值误差则直接改变被聚合的内容。两者经过后续层和未来生成反馈还可能继续传播，所以不能只看两组张量的平均重建误差就认定对模型影响相同。

量化粒度可以按张量、头、通道、词元或小组划分。更细粒度通常更容易适应局部数值范围，却需要更多元数据，并使内核更复杂。KIVI 研究了键和值的分布差异，采用键侧按通道、值侧按词元的非对称设计，并在其测试条件下评估低位宽缓存\cite{kivi}。这是一项有具体对象与实现的研究结果，不是所有模型、所有位置编码和所有任务都应固定使用同一位宽的普遍保证。

\subsection{何时有望加速，何时只是节省容量}

若解码主要受缓存读取带宽限制，低位宽缓存有机会同时降低显存和读取字节数；若内核每步先把全部历史反量化为高精度临时张量再计算，则额外读写可能吞掉收益。短序列、小批量时，量化和反量化的固定成本也可能超过节省的访问时间。高效实现通常需要让解码计算与压缩布局协同，而不是在任意缓存对象外面加一次类型转换。

质量评估应同时考察短文本、长距离信息检索、多证据推理与较长输出，比较量化前后固定输入上的分数差异以及自由生成后的任务表现。只比较几条看起来通顺的回答，不能证明缓存量化无损；只比较困惑度也未必能覆盖特定业务中对少量关键词元的敏感性。工程决策应给定可接受的质量变化、显存收益和延迟要求，并在这些条件下选择位宽与分组大小。

\section{长上下文中的窗口、淘汰与换出}
\label{sec:kv-long-context}

\subsection{原生局部注意力与事后截断}

\keyterm{滑动窗口注意力（sliding-window attention）}让某层当前位置只读取最近的 $W$ 个位置。若模型本来就按这一依赖关系定义，过去已经离开窗口、且未来再也不会访问的位置可以释放，该层的长期缓存规模被窗口限制。若一个模型原本使用全局注意力，却在推理时临时丢弃早期键值，则改变了允许访问的集合，通常不再与原模型等价。这两种情况都可能叫“窗口缓存”，但一个是在利用原模型的依赖结构，另一个是在近似原模型。

对于多层局部注意力，较新的高层状态可能间接汇集更早的信息，因此可直接读取的窗口不等于整个网络的有效感受范围。不过，间接编码也不保证完整保存所有早期细节，更不能推出任意远处信息都可精确恢复。混合局部层与全局层时，各层保留长度不同，应使用式\eqref{eq:kv-memory-heterogeneous}，不能把最短窗口直接乘到全部层上。

事后裁剪还与“把保留后缀当作新提示重新计算”不同：保留下来的高层键值在最初计算时可能读取过后来被删除的前文，重新计算后缀则没有这些依赖。两种状态可能得到不同输出。应用必须明确自己的目标是保留原运行轨迹的部分状态，还是重新定义一个截断后的输入问题；两者都可用于工程，但不能未经证明相互替换。

\subsection{重要性淘汰与信息不可逆性}

\keyterm{缓存淘汰（cache eviction）}按某种准则删除部分键值，例如只保留最近位置，或结合历史注意力强度保留一些经常被读取的位置。H$_2$O 研究了结合近期位置与累计重要位置的淘汰策略\cite{h2o}；StreamingLLM 则研究了保留少量初始位置及近期窗口所产生的流式生成行为\cite{streaming}。它们针对特定模型和任务展示了容量与质量之间的权衡，不等价于在固定内存里无损保存无限历史。

\keyterm{注意力汇聚位置（attention sink）}指某些位置可能获得较强注意力，即使其文字内容未必是当前任务的重要证据。保留这些位置有时有助于稳定模型行为，但“稳定地继续生成”与“能够回答任意远处的细节问题”是不同目标。若一个早期随机编号已从所有可访问状态中丢失，后续再询问它时，模型不能靠通顺语言证明自己仍记得编号。评价流式方案应包含对旧证据的实际访问测试，而不只记录能生成多长文本。

过去的重要性也不一定预测未来需求。一段此前几乎没有被注意的合同条款，可能在后来的具体问题中突然成为唯一证据；一个过去高频出现的模板词，则可能对后续答案不再有用。在线淘汰不能预知所有未来查询，因此任何经验准则都应说明保留预算、任务假设和失败模式。若任务要求对任意历史事实进行精确查证，保留原始文本并在需要时检索或重算，通常比把有损缓存当作完整档案更符合该目标。

\subsection{换出与重新计算的代价}

\keyterm{缓存换出（cache offloading）}把暂不驻留设备的状态放到主机内存或其他存储，在需要时再传入。如果字节内容完整保留、读取顺序与数值计算不变，换出本身不引入量化或删除误差；代价主要是传输与调度。若一次必须传输 $D_{\mathrm{link}}$ 字节，链路有效带宽为 $\beta_{\mathrm{link}}$，仅数据搬运就至少需要约 $D_{\mathrm{link}}/\beta_{\mathrm{link}}$ 的时间，另有启动与同步成本。能否与计算重叠，要看预取顺序、缓冲空间和真实并行能力。

按层预取可以让当前层计算与下一层数据传输部分重叠，但无法凭空消除带宽需求。长上下文每步都读取大量历史时，设备与主机之间的链路可能成为新的瓶颈。因此，换出首先扩展的是可运行容量，不自动提升单请求速度。它对于允许更高延迟、原本无法装入的请求仍可能有价值；讨论收益时应比较同等可完成任务的方案，而不只比较理想的全驻留情形。

丢弃缓存后从词元重新计算是另一种时间换空间方式。原始词元通常远小于各层键值，但恢复计算可能需要从更早的有效前缀开始，并不能只根据被删位置的词元编号立即恢复该位置高层键值。选择换出还是重算，应比较传输时间、提示重算成本、设备拥塞与请求优先级；不应把缓存字节数小的方案自动视为总成本最低。

\section{分支生成、分布式推理与训练边界}
\label{sec:kv-extensions}

\subsection{束搜索与推测解码的状态回滚}

\keyterm{束搜索（beam search）}同时维持多个候选前缀。相同父前缀可以共享只读缓存，但每个分支追加不同词元后会产生独立状态。每轮选择新候选时，需要按父候选索引重排缓存，且一个父候选可能被复制为多个子分支。仅重排累计分数或生成词元，而不重排缓存，会把一条分支的历史接到另一条分支上。共享块和写时复制可以减少公共部分的重复保存，却不能消除所有分支后缀的成本。

\keyterm{推测解码（speculative decoding）}先由较便宜的机制提出多个候选词元，再让目标模型成批验证。特定接受与修正算法可以保持目标模型的采样分布，其条件与证明见文献\cite{speculative}；任意接受草稿都不具有这种保证。对缓存而言，验证时可能临时写入一串候选状态，最终只接受其中一个前缀，因此必须撤销未接受部分，并与真正输出的词元边界对齐。

例如有效历史长度为 $T$，验证三个候选后只接受前两个，那么目标模型缓存中超过已接受且已实际处理边界的部分不能继续作为历史。若随后选择了一个修正词元，它与被拒绝的第三个候选不同，应在正确前缀上重新处理。具体实现可能暂时少缓存一个已输出词元，与第\ref{sec:kv-phases}节相同；关键是显式记录已接受、已输出和已处理三个计数。草稿模型和目标模型参数不同，通常各自维护自己的缓存，不能只因词元序列一致就直接互换键值。

\subsection{多设备分片不是简单除以设备数}

\keyterm{张量并行（tensor parallelism）}可以把头或特征维分到不同设备，\keyterm{流水线并行（pipeline parallelism）}可以把不同层放在不同设备。若键值头在 $p$ 台设备上恰好均匀切分且没有复制，每台设备的相应缓存载荷可以约为总量的 $1/p$；若 $H_{kv}$ 很少、无法按设备数均分，或实现复制共享键值头，则实际缩减比例不同。流水线划分则应按每台设备拥有的层数求和。容量估算必须读取真实分片方式，不能直接把总显存相加后认为任意一层都能使用全部容量。

按上下文位置分片还可能要求跨设备交换注意力统计量或数据。式\eqref{eq:kv-online-softmax}给出了组合不同键段结果的数学基础，但分布式归约会增加通信和同步。缓存被分散保存后，局部容量压力减小，端到端延迟却取决于网络以及负载均衡。单机显存公式仍有用，但它只是通信分析的起点。

把预填充与解码放在不同设备池，能够针对两个阶段的计算特征分配资源，但提示处理结束后需要转移对应状态。转移量与提示缓存规模有关，长提示的状态传输可能抵消调度收益。若公共前缀在解码侧已经存在，可以研究减少重复传输，但必须验证身份与格式兼容。预填充和解码分离是一项服务架构选择，不是打开 KV-cache 后自然获得的属性。

\subsection{交叉注意力与训练激活}

编码器--解码器模型的\keyterm{交叉注意力（cross-attention）}中，查询来自不断增长的解码器，键和值来自固定的编码器输出。源输入与参数不变时，各解码层对应的交叉注意力键值可以预先计算并复用，其长度取决于源序列，不随目标生成同步增长；解码器自注意力缓存则仍随目标前缀增长。两种缓存的来源和长度不同，计算总内存时应分别求和。

训练中常使用\keyterm{教师强制（teacher forcing）}，一次输入已知目标序列的移位版本，并行计算各位置损失，因此没有与逐词元生成相同的重复历史前向。反向传播还需要相应计算图和梯度关系；将历史键值直接分离并跨训练步长期保存，可能切断本应存在的梯度，参数更新后也会让旧状态失效。特殊的分块训练、记忆机制或冻结前缀方案需要单独定义目标与梯度边界，不能把推理缓存无条件搬进普通训练循环。

因此，\keyterm{激活检查点（activation checkpointing）}与 KV-cache 的方向不同：前者在训练中少存一些中间激活并在反向时重算，以节省内存；后者在推理中多存可复用的历史键值，以减少未来重算。两者都在时间与空间之间取舍，但适用阶段、保存对象和正确性条件各不相同。

\section{实验设计、诊断方法与结果解释}
\label{sec:kv-validation}

\subsection{先验证固定轨迹上的函数一致性}

最小正确性实验应先冻结一段输入词元序列，而不是让两个实现自由采样。对每一个前缀，基线从头完整计算，待测实现复用已有缓存并只处理新增位置，逐步比较最后位置的分数、每层键值和长度。这样不会因早期一次采样分叉而把后续不同输入上的输出误差当成数值误差。除逐词元追加外，还要比较一次性输入与不同分块方案，特别是历史长度不为零且新块长度大于一的矩形掩码情形。

设基线分数为 $\mathbf z$，待测分数为 $\widehat{\mathbf z}$，可记录最大绝对误差 $\max_i|z_i-\widehat z_i|$，以及逐元素容差条件 $|z_i-\widehat z_i|\leq\varepsilon_{\mathrm{abs}}+\varepsilon_{\mathrm{rel}}|z_i|$ 的通过情况，其中两种容差均须预先给定并与数据精度匹配。验证还应覆盖空缓存、单位置提示、达到容量边界、请求结束后重用槽位以及不同长度批次。通过一个短例子，只证明该例子的路径可用，不构成对所有配置的证明。

若只有第一个新词元之后开始错误，优先检查位置偏移、末尾词元是否重复处理以及缓存长度；若单词元追加正确而多词元追加错误，优先检查矩形因果掩码；若单请求正确而批量错误，检查填充、有效长度与请求映射；若分支生成才错误，检查父索引、共享写入与回滚。此类定位依据计算依赖，而不是根据最终文字是否“像正常回答”判断。

对数值近似方案，应把两种实验分开：固定词元轨迹用于测量数值变化，自由生成用于评估变化经过采样反馈后的任务质量。关闭梯度计算与关闭训练随机性也应分别确认：不构建梯度图不必然自动切换模型的训练模式。比较时应固定模型版本、分词器、模板、位置配置、内核和精度；改变这些因素会使实验同时比较多个变量。

\subsection{性能指标必须对应用户等待过程}

\keyterm{首词元延迟（time to first token, TTFT）}是从规定的请求起点到首次输出词元的时间，端到端口径可能包含网络、排队、分词和预填充；内核口径则可能只算设备执行，两者应明确区分。\keyterm{词元间延迟（inter-token latency, ITL）}描述相邻输出词元之间的等待；\keyterm{每输出词元时间（time per output token, TPOT）}常用首词元之后的生成耗时除以剩余输出数概括平均速度。平均值不能反映停顿，因此在线服务还应报告尾部分位数。

\keyterm{吞吐（throughput）}需要注明分子是输入词元、输出词元还是请求数，分母是设备忙时还是整个观测窗口。批量增加可能提高总输出词元每秒，却同时增加某条请求的首词元等待与相邻词元间隔。优化目标通常应是满足延迟约束下的可承载吞吐，而不是单一峰值。若实验固定总输出长度，却忽略缓存开启后可以提高并发的空间，也可能低估系统收益；反之，只比较不同并发下的峰值，又无法隔离单请求内核变化。

异步设备执行使普通主机计时容易只测到任务提交时间。内核实验需要适当同步或设备事件，先完成必要预热，再测稳定阶段；服务实验还应记录到达过程和排队。缓存、图编译和前缀复用的预热状态应公开，冷启动与热启动分别报告。提前终止、输出长度变化、采样随机性和后台任务都可能影响结果，不能由一次计时得出普遍加速比例。

\subsection{一个可执行的对照矩阵}

可以先选定同一模型与精度，在 $P\in\{128,1024,8192\}$、$G\in\{1,32,128\}$ 和若干可容纳批量上比较缓存开关；这些长度是实验设计示例，应按模型上下文范围与设备容量调整。$G=1$ 检查只做预填充的边界，较大 $G$ 显示历史复用收益，长 $P$ 揭示缓存容量和读取压力。每项记录峰值显存、预填充耗时、解码耗时和正确性误差；超出容量的情况记为无法完成，不能悄悄缩短输入后仍放在同一行比较。

然后在已验证的缓存基线上逐项加入静态分配、分页、前缀共享、量化或换出，每次只改变一个主因素，并保留必要的组合实验。前缀共享同时测试无共享、全部共享和部分共享；量化同时测试质量与速度；淘汰需要额外测试远处证据；分页应测不同长度分布与请求进入退出。这样每项结果才能说明该机制在什么条件下有效，而不是仅得出“某个复杂配置更快”的不可解释结论。

最后根据观测选择改进方向：历史重算占主导时先采用正确缓存；长期载荷过大时审视键值头、精度与保留策略；分配浪费较大时考虑分页和共享；传输占主导时减少不必要换出与跨设备迁移；数值或质量下降时回到掩码、位置和有损近似的控制实验。优化应由已测瓶颈驱动，不必把所有方法同时打开。本文未给出任何硬件实测速率，前面的数值都是明确假设下的解析计算或教学示例。

\section{综合练习与参考解答}
\label{sec:kv-exercises}

\begin{example}[由容量倒推并发预算]
某模型有 $L=40$ 层、$H_{kv}=8$、$d_h=128$，键和值均为每元素两字节。缓存预算为 $10\,\mathrm{GiB}$，请求间没有共享，每条请求峰值保留 $4096$ 个词元。忽略块表与对齐，最多容纳多少条这样的请求？
\end{example}
\begin{solve}
每词元缓存为 $2\times40\times8\times128\times2=163840$ 字节，即 $160\,\mathrm{KiB}$。每请求需要 $640\,\mathrm{MiB}$，而 $10\,\mathrm{GiB}=10240\,\mathrm{MiB}$，因此载荷上限为十六条。若给出的预算尚未扣除工作区或元数据，则实际可接纳数量应进一步减少。解题时使用的是键值头数，查询头数不出现在这一步载荷计数中。
\end{solve}

\begin{example}[前缀共享节约了什么]
四条请求均有 $2048$ 个词元的公共前缀，各自已处理 $512$ 个不同后缀词元；每词元缓存为 $128\,\mathrm{KiB}$，假定公共部分恰能完整共享。求共享前后的载荷，并解释是否可据此推出四倍解码速度。
\end{example}
\begin{solve}
不共享需保存 $4(2048+512)=10240$ 份词元状态，占 $1280\,\mathrm{MiB}$。共享后保存 $2048+4\times512=4096$ 份，占 $512\,\mathrm{MiB}$，节约 $768\,\mathrm{MiB}$。每条请求的新查询仍然要读取其可见前缀并计算自己的注意力权重，因此不能推出四倍解码速度；容量节约和预填充复用可以改善系统吞吐，但需要实际调度与内核实验支持。
\end{solve}

\begin{example}[判断三种复用是否合法]
同一个模型处理两个输入：其一具有相同词元前缀但采用不同采样温度；其二文本字符串相同但聊天模板不同；其三前缀词元相同但切换了影响注意力投影的适配器。哪些前缀缓存可以直接复用？
\end{example}
\begin{solve}
第一种在其他前向条件相同且前缀固定时可以复用，因为普通温度调整作用于选择下一词元时的分数，不改变该固定前缀的键值；以后若采样结果不同，分叉之后的状态各自生成。第二种必须比较模板处理后的实际输入，字符串相同不能保证词元前缀和角色控制信息相同。第三种改变了产生键值的函数，不能直接使用原适配器的缓存。判断依据是计算是否相同，而不是请求文本或会话身份是否相似。
\end{solve}

\begin{example}[设计一个能发现掩码错误的测试]
已有三个历史位置，一次追加两个位置。某实现声称已开启因果模式，却使用矩形矩阵左上角对齐的下三角掩码。说明一个无需训练大模型的检验方法。
\end{example}
\begin{solve}
令全部查询键分数为零，五个值依次为 $1,2,3,4,5$。正确的两行输出分别为前四个值的均值 $2.5$ 和前五个值的均值 $3$；错误掩码则输出 $1$ 和 $1.5$。这个小例子不依赖模型学到什么，只检查可见集合和归一化，因而比观察一段自然语言输出更容易定位掩码对齐问题。通过后仍需测试非零分数、不同头和实际位置编码。
\end{solve}

\begin{example}[解释缓存没有解决的三件事]
为什么开启 KV-cache 之后，生成仍然串行、长上下文仍然变慢，而且模型仍然可能忘记提示中的事实？
\end{example}
\begin{solve}
标准自回归生成的下一个输入依赖上一步选择的词元，缓存没有取消这一依赖；全局注意力中新查询仍要读取增长的历史，缓存没有取消每步的历史访问；模型是否正确利用某条事实取决于训练、表示和任务推理，保存键值只保证相应状态可供计算，并不保证模型把注意力或答案分配给正确证据。计算复用、资源效率和任务正确性因此应分别验证。
\end{solve}

\section{文献范围与进一步阅读}

学习顺序宜先掌握式\eqref{eq:kv-attention-row}、命题\ref{prop:kv-prefix-invariant}与生成循环，再利用式\eqref{eq:kv-memory-general}估算实际模型容量，最后进入系统优化。论文用于核对提出的方法与适用条件，官方文档用于核对软件对象和接口。本文对论文方法的介绍仅说明与缓存相关的核心机制，未复制其完整算法证明或把原论文实验结果推广到任意模型。软件文档链接按 2026 年 10 月 3 日检索，滚动更新页面的接口和支持范围可能随后变化，实际复现应固定软件版本。

\begin{thebibliography}{99}
\bibitem{transformer} Vaswani A., et al. Attention Is All You Need. 2017. \url{https://arxiv.org/abs/1706.03762}.
\bibitem{hf-cache} Hugging Face. How caching works. Transformers 官方文档. \url{https://huggingface.co/docs/transformers/main/en/cache_explanation}.
\bibitem{hf-strategies} Hugging Face. Cache strategies. Transformers 官方文档. \url{https://huggingface.co/docs/transformers/main/en/kv_cache}.
\bibitem{mqa} Shazeer N. Fast Transformer Decoding: One Write-Head is All You Need. 2019. \url{https://arxiv.org/abs/1911.02150}.
\bibitem{gqa} Ainslie J., et al. GQA: Training Generalized Multi-Query Transformer Models from Multi-Head Checkpoints. 2023. \url{https://arxiv.org/abs/2305.13245}.
\bibitem{rope} Su J., et al. RoFormer: Enhanced Transformer with Rotary Position Embedding. 2021 年预印本及后续修订. \url{https://arxiv.org/abs/2104.09864}.
\bibitem{flash} Dao T., et al. FlashAttention: Fast and Memory-Efficient Exact Attention with IO-Awareness. 2022. \url{https://arxiv.org/abs/2205.14135}.
\bibitem{paged} Kwon W., et al. Efficient Memory Management for Large Language Model Serving with PagedAttention. 2023. \url{https://arxiv.org/abs/2309.06180}.
\bibitem{vllm-prefix} vLLM. Automatic Prefix Caching. 官方设计文档. \url{https://docs.vllm.ai/en/latest/design/prefix_caching/}.
\bibitem{mla} DeepSeek-AI. DeepSeek-V2: A Strong, Economical, and Efficient Mixture-of-Experts Language Model. 2024. \url{https://arxiv.org/abs/2405.04434}.
\bibitem{kivi} Liu Z., et al. KIVI: A Tuning-Free Asymmetric 2bit Quantization for KV Cache. 2024. \url{https://arxiv.org/abs/2402.02750}.
\bibitem{h2o} Zhang Z., et al. H$_2$O: Heavy-Hitter Oracle for Efficient Generative Inference of Large Language Models. 2023. \url{https://arxiv.org/abs/2306.14048}.
\bibitem{streaming} Xiao G., et al. Efficient Streaming Language Models with Attention Sinks. 2023 年预印本及后续修订. \url{https://arxiv.org/abs/2309.17453}.
\bibitem{speculative} Leviathan Y., Kalman M., Matias Y. Fast Inference from Transformers via Speculative Decoding. 2023. \url{https://arxiv.org/abs/2211.17192}.
\end{thebibliography}
\end{document}
